\chapter{Menciptakan Makna Bersama: Shared Meaning System}

\begin{insightbox}
\textit{``Hubungan yang benar-benar sukses bukan hanya tentang bertahan dari konflik atau membangun keintiman harian. Puncaknya adalah ketika kalian menciptakan \highlight{makna bersama}---sebuah budaya hubungan yang unik, ritual yang mengikat, dan tujuan hidup yang dibangun berdua. Ini adalah Prinsip 7: cita-cita tertinggi dari hubungan yang berkembang.''} --- Dr. John Gottman
\end{insightbox}

\section{Apa itu Shared Meaning?}

\keyterm{Shared meaning} adalah sistem nilai, simbol, dan tujuan bersama yang menciptakan identitas unik bagi sebuah hubungan. Ini melampaui sekadar menjalani hari demi hari---ini tentang \highlight{kehidupan bersama dengan tujuan}.

\begin{praktikbox}[Definisi Shared Meaning yang Mendalam]
Shared meaning terdiri dari empat pilar utama:

\begin{enumerate}
    \item \textbf{Rituals of Connection:} Tradisi dan rutinitas yang mengikat
    \item \textbf{Roles:} Makna dari menjadi ``suami/istri/pasangan''
    \item \textbf{Goals:} Impian dan tujuan hidup yang dibangun bersama
    \item \textbf{Symbols:} Arti dari rumah, uang, waktu, dan objek bersama
\end{enumerate}

\textit{Shared meaning bukan sekadar kesamaan---ini tentang saling menghormati perbedaan sambil menciptakan kehidupan yang bermakna bersama.}
\end{praktikbox}

\subsection{Shared Meaning vs Daily Logistics}

Banyak pasangan mengira mereka punya shared meaning hanya karena mereka:
\begin{itemize}
    \item Hidup di rumah yang sama
    \item Membagi tugas rumah tangga
    \item Mengurus anak bersama
    \item Membayar tagihan bersama
\end{itemize}

\begin{warningbox}[Perbedaan Fundamental]
\textbf{Daily Logistics (Operasional):}
\begin{itemize}
    \item ``Siapa yang ambil anak hari ini?''
    \item ``Berapa budget bulan ini?''
    \item ``Acara keluarga minggu depan kita pergi?''
\end{itemize}

\textbf{Shared Meaning (Filosofis):}
\begin{itemize}
    \item ``Apa makna menjadi orang tua bagi kita?''
    \item ``Bagaimana uang mencerminkan nilai kita?''
    \item ``Tradisi keluarga apa yang ingin kita bangun?''
\end{itemize}

\textit{Pasangan dengan shared meaning kuat menjalani \textbf{kehidupan yang sama}. Pasangan tanpa shared meaning hanya \textbf{hidup berdampingan}.}
\end{warningbox}

\section{The Culture of Relationship: Membangun Budaya Hubungan}

Setiap hubungan memiliki \keyterm{budaya}---serangkaian aturan tak tertulis, nilai, dan praktik yang mengatur bagaimana pasangan hidup bersama. Budaya ini bisa dibangun dengan sengaja atau terbentuk secara tak sadar.

\begin{center}
\begin{tikzpicture}[
    culture/.style={draw=#1, line width=2.5pt, rounded corners=15pt,
                    fill=#1!10, minimum width=13cm, minimum height=2cm,
                    align=left, text width=12cm}
]
    \node[culture=primary] (ritual) at (0,6) {
        \textbf{\large \textcolor{primary}{1. RITUALS OF CONNECTION}}\\[0.2cm]
        \small Tradisi, rutinitas, dan momen-momen sakral yang mengikat
    };
    
    \node[culture=secondary] (roles) at (0,3.5) {
        \textbf{\large \textcolor{secondary}{2. ROLES}}\\[0.2cm]
        \small Definisi dan ekspektasi dari peran masing-masing dalam hubungan
    };
    
    \node[culture=accent] (goals) at (0,1) {
        \textbf{\large \textcolor{accent}{3. GOALS}}\\[0.2cm]
        \small Impian bersama, visi jangka panjang, dan legacy yang dibangun
    };
    
    \node[culture=sagegreen] (symbols) at (0,-1.5) {
        \textbf{\large \textcolor{sagegreen}{4. SYMBOLS}}\\[0.2cm]
        \small Makna dari tempat, objek, dan konsep dalam hubungan
    };
    
    \node[draw=warmgray, line width=3pt, rounded corners=15pt, 
          fill=warmgray!10, minimum width=10cm, minimum height=1.5cm, align=center] 
          at (0,-4) {\textbf{\large SHARED MEANING SYSTEM}};
    
    \draw[->, line width=2pt] (ritual.south) -- (roles.north);
    \draw[->, line width=2pt] (roles.south) -- (goals.north);
    \draw[->, line width=2pt] (goals.south) -- (symbols.north);
    \draw[->, line width=2pt] (symbols.south) -- ++(0,-0.7);
\end{tikzpicture}
\end{center}

\subsection{Mengapa Budaya Hubungan Penting?}

\begin{praktikbox}[Benefit Shared Meaning System]
\textbf{1. Sense of Unity}
\begin{itemize}
    \item ``Kita'' menjadi lebih penting dari ``Aku'' vs ``Kamu''
    \item Identitas sebagai tim terbentuk
    \item Loyalitas pada hubungan meningkat
\end{itemize}

\textbf{2. Resilience During Stress}
\begin{itemize}
    \item Ritual menjadi anchor saat krisis
    \item Nilai bersama memberi arah dalam keputusan sulit
    \item Shared meaning sebagai ``why'' yang memotivasi bertahan
\end{itemize}

\textbf{3. Protection Against Drift}
\begin{itemize}
    \item Mencegah hubungan menjadi ``roommate arrangement''
    \item Menjaga romance di tengah rutinitas
    \item Memberi makna pada usaha-usaha kecil sehari-hari
\end{itemize}
\end{praktikbox}

\section{Self-Assessment: Seberapa Kuat Shared Meaning-mu?}

\begin{exercisebox}[Shared Meaning Assessment]
\textbf{Petunjuk:} Jawab setiap pertanyaan dengan skala 1-5 (1 = Sangat tidak setuju, 5 = Sangat setuju)

\subsection*{BAGIAN A: Rituals (1-5 poin per item)}
\begin{enumerate}
    \item Kami memiliki rutinitas pagi yang konsisten bersama
    \item Kami punya tradisi akhir pekan yang kami nantikan
    \item Ada ritual khusus saat kami berpisah dan bertemu kembali
    \item Kami merayakan anniversary dengan cara yang bermakna
    \item Ada ritual sebelum tidur yang menguatkan koneksi kami
    \item Kami punya tradisi liburan tahunan yang konsisten
\end{enumerate}

\subsection*{BAGIAN B: Roles (1-5 poin per item)}
\begin{enumerate}
    \item[7.] Kami sepakat tentang makna ``suami'' dan ``istri'' dalam hubungan kami
    \item[8.] Peran kami sebagai orang tua (jika ada anak) sejalan dengan nilai kami
    \item[9.] Kami jelas tentang ekspektasi satu sama lain dalam rumah tangga
    \item[10.] Kami merasa bangga dengan peran kami masing-masing
    \item[11.] Ada pembagian tanggung jawab yang fair dan disepakati
\end{enumerate}

\subsection*{BAGIAN C: Goals (1-5 poin per item)}
\begin{enumerate}
    \item[12.] Kami punya visi jangka panjang yang sama tentang masa depan
    \item[13.] Kami sering mendiskusikan impian dan aspirasi hidup
    \item[14.] Tujuan karir/finansial kami saling mendukung
    \item[15.] Kami punya proyek bersama yang sedang dikerjakan
    \item[16.] Legacy apa yang ingin kami tinggalkan sudah jelas
\end{enumerate}

\subsection*{BAGIAN D: Symbols (1-5 poin per item)}
\begin{enumerate}
    \item[17.] Rumah kami mencerminkan nilai dan identitas kami berdua
    \item[18.] Kami setuju tentang makna uang dalam hidup kami
    \item[19.] Waktu bersama dihargai sebagai komoditas yang berharga
    \item[20.] Ada objek/ruang khusus yang punya makna simbolis bagi kami
\end{enumerate}

\vspace{0.5cm}
\textbf{SCORING:}

\begin{center}
\renewcommand{\arraystretch}{1.4}
\begin{tabular}{|c|c|p{7cm}|}
\hline
\rowcolor{primary!20}
\textbf{Skor Total} & \textbf{Level} & \textbf{Interpretasi} \\
\hline
80-100 & \textbf{Sangat Kuat} & Shared meaning system sangat terbentuk. Hubungan punya fondasi filosofis yang kokoh \\
\hline
60-79 & \textbf{Kuat} & Ada fondasi yang baik, tapi masih ada area untuk ditingkatkan \\
\hline
40-59 & \textbf{Sedang} & Shared meaning mulai terbentuk tapi belum solid. Perlu sengaja membangun \\
\hline
20-39 & \textbf{Lemah} & Shared meaning minim. Risiko drift dan kehampaan hubungan \\
\hline
\textless 20 & \textbf{Sangat Lemah} & Shared meaning hampir tidak ada. Intervensi segera diperlukan \\
\hline
\end{tabular}
\end{center}
\end{exercisebox}

\section{Level 1: Rituals of Connection}

\keyterm{Rituals} adalah tindakan yang diulang dengan makna simbolis---bukan sekadar kebiasaan, tapi praktik yang sengaja dibangun untuk memperkuat ikatan.

\subsection{Jenis-Jenis Ritual}

\begin{center}
\renewcommand{\arraystretch}{1.5}
\begin{tabular}{|>{\raggedright}p{2.5cm}|>{\raggedright}p{4.5cm}|>{\raggedright}p{4.5cm}|>{\raggedright\arraybackslash}p{3.5cm}|}
\hline
\rowcolor{primary!20}
\textbf{Frekuensi} & \textbf{Contoh Ritual} & \textbf{Fungsi} & \textbf{Pentingnya} \\
\hline
\textbf{Daily} & Morning coffee together, Goodbye kiss, Dinner ritual, Bedtime routine & Menjaga koneksi harian & \textcolor{sagegreen}{Sangat Tinggi} \\
\hline
\textbf{Weekly} & Date night, Sunday brunch, Weekend walk & Recharge dan quality time & \textcolor{sagegreen}{Tinggi} \\
\hline
\textbf{Monthly} & Relationship check-in, Adventure day, Financial review & Maintenance dan planning & \textcolor{secondary}{Sedang} \\
\hline
\textbf{Annual} & Anniversary celebration, Vacation ritual, Year-end reflection & Milestone dan renewal & \textcolor{secondary}{Sedang} \\
\hline
\textbf{Life Events} & Birth rituals, Career celebration, Grieving together & Transisi hidup & \textcolor{sagegreen}{Sangat Tinggi} \\
\hline
\end{tabular}
\end{center}

\subsection{Membangun Ritual dari Nol}

\begin{tipbox}[Panduan Membangun Ritual Baru]
\textbf{Step 1: Identifikasi Gap}
\begin{itemize}
    \item Kapan kalian merasa paling terhubung?
    \item Kapan kalian merasa paling terpisah?
    \item Apa ritual dari masa kecil yang ingin kalian lanjutkan atau hindari?
\end{itemize}

\textbf{Step 2: Desain Ritual}
\begin{itemize}
    \item Pilih waktu yang konsisten
    \item Tentukan struktur (apa yang terjadi, dalam urutan apa)
    \item Buat elemen simbolis (objek, lagu, frasa khusus)
\end{itemize}

\textbf{Step 3: Implementasi}
\begin{itemize}
    \item Mulai sederhana---jangan terlalu kompleks
    \item Konsisten selama 21 hari untuk membentuk kebiasaan
    \item Evaluasi dan adjust setelah sebulan
\end{itemize}

\textbf{Step 4: Proteksi}
\begin{itemize}
    \item Jadikan ritual ``sacred''---prioritas non-negotiable
    \item Komunikasikan pada orang lain (anak, kerjaan) bahwa ini penting
    \item Jaga ritual bahkan saat stres atau konflik
\end{itemize}
\end{tipbox}

\subsection{Template Desain Ritual}

\begin{exercisebox}[Worksheet: Membuat Ritual Baru]
\textbf{Nama Ritual:} \rule{8cm}{0.4pt}

\textbf{Frekuensi:} $\square$ Daily $\square$ Weekly $\square$ Monthly $\square$ Annual $\square$ Special occasion

\textbf{Waktu/Trigger:} \rule{8cm}{0.4pt}

\textbf{Struktur Ritual:}
\begin{enumerate}
    \item \rule{12cm}{0.4pt}
    \item \rule{12cm}{0.4pt}
    \item \rule{12cm}{0.4pt}
\end{enumerate}

\textbf{Elemen Simbolis:} (lagu, objek, frasa, gesture)
\begin{itemize}
    \item \rule{12cm}{0.4pt}
    \item \rule{12cm}{0.4pt}
\end{itemize}

\textbf{Makna Ritual Ini:} (mengapa ini penting)
\begin{itemize}
    \item \rule{12cm}{0.4pt}
\end{itemize}

\textbf{Proteksi:} Bagaimana menjaga ritual ini?
\begin{itemize}
    \item \rule{12cm}{0.4pt}
\end{itemize}

\textbf{Ditandatangani:}\\
Partner A: \rule{5cm}{0.4pt} Partner B: \rule{5cm}{0.4pt}
\end{exercisebox}

\section{Level 2: Roles---Makna dari Menjadi Pasangan}

\keyterm{Roles} dalam shared meaning adalah definisi tentang apa artinya menjadi ``suami,'' ``istri,'' atau ``pasangan'' dalam hubungan spesifik kalian. Ini tentang ekspektasi, tanggung jawab, dan identitas.

\subsection{Clarifying Roles}

Setiap orang membawa ekspektasi tentang peran dari:
\begin{itemize}
    \item \textbf{Keluarga asal:} Bagaimana orang tua mereka menjalankan peran
    \item \textbf{Budaya:} Norma gender dan peran dalam masyarakat
    \item \textbf{Pengalaman pribadi:} Hubungan sebelumnya dan pertumbuhan personal
    \item \textbf{Media:} Representasi peran dalam film, buku, media sosial
\end{itemize}

\begin{praktikbox}[Eksplorasi Peran: Pertanyaan Kunci]
\textbf{Makna ``Menjadi Suami/Partner A'':}
\begin{enumerate}
    \item Apa tanggung jawab utama sebagai suami menurutmu?
    \item Bagaimana suami yang ideal menunjukkan cinta?
    \item Apa peran suami dalam pengambilan keputusan?
    \item Bagaimana suami menyeimbangkan kerja dan keluarga?
    \item Apa yang membuat seorang suami merasa berhasil?
\end{enumerate}

\textbf{Makna ``Menjadi Istri/Partner B'':}
\begin{enumerate}
    \item Apa tanggung jawab utama sebagai istri menurutmu?
    \item Bagaimana istri yang ideal menunjukkan cinta?
    \item Apa peran istri dalam pengambilan keputusan?
    \item Bagaimana istri menyeimbangkan karir dan keluarga?
    \item Apa yang membuat seorang istri merasa berhasil?
\end{enumerate}
\end{praktikbox}

\subsection{Roles sebagai Orang Tua (Jika Berlaku)}

\begin{center}
\renewcommand{\arraystretch}{1.5}
\begin{tabular}{|>{\raggedright}p{3.5cm}|>{\raggedright}p{5cm}|>{\raggedright\arraybackslash}p{6cm}|}
\hline
\rowcolor{sagegreen!20}
\textbf{Aspek Parenting} & \textbf{Diskusi yang Perlu Dilakukan} & \textbf{Tujuan} \\
\hline
\textbf{Philosophy} & Apa tujuan utama parenting kami? & Memastikan arah yang sama \\
\hline
\textbf{Discipline} & Gaya disiplin seperti apa? Kapan intervensi? & Konsistensi dalam parenting \\
\hline
\textbf{Nurturing} & Bagaimana menunjukkan kasih sayang? & Memenuhi kebutuhan emosional anak \\
\hline
\textbf{Education} & Visi pendidikan dan perkembangan anak & Investasi jangka panjang \\
\hline
\textbf{Division of Labor} & Siapa bertanggung jawab apa? & Fairness dan teamwork \\
\hline
\end{tabular}
\end{center}

\section{Level 3: Goals---Impian dan Tujuan Bersama}

\keyterm{Goals} dalam shared meaning adalah impian hidup yang dibangun bersama---bukan sekadar target pribadi, tapi visi yang diselaraskan dan dikejar sebagai tim.

\subsection{Hierarki Tujuan dalam Hubungan}

\begin{center}
\begin{tikzpicture}[
    goal/.style={draw=#1, line width=2pt, rounded corners=10pt, fill=#1!10,
                 minimum width=10cm, minimum height=1.5cm, align=center}
]
    \node[goal=primary] (legacy) at (0,6) {
        \textbf{LEGACY GOALS}\\
        \small Warisan, nilai yang ditinggalkan, kontribusi pada dunia
    };
    
    \node[goal=secondary] (life) at (0,4) {
        \textbf{LIFE DREAMS}\\
        \small Impian besar: karir, pendidikan, pengalaman hidup
    };
    
    \node[goal=accent] (shared) at (0,2) {
        \textbf{SHARED PROJECTS}\\
        \small Proyek bersama: rumah, bisnis, komunitas
    };
    
    \node[goal=sagegreen] (annual) at (0,0) {
        \textbf{ANNUAL GOALS}\\
        \small Target tahunan yang konkret dan terukur
    };
    
    \draw[->, line width=2pt] (annual) -- (shared);
    \draw[->, line width=2pt] (shared) -- (life);
    \draw[->, line width=2pt] (life) -- (legacy);
\end{tikzpicture}
\end{center}

\subsection{Shared Goals Exercise}

\begin{exercisebox}[Worksheet: Goals Alignment]
\textbf{Petunjuk:} Jawab secara individual, lalu diskusikan dan cari overlap.

\subsection*{A. Personal Dreams (Individual)}
\begin{enumerate}
    \item Jika uang dan waktu tidak masalah, apa yang ingin kamu capai dalam 5 tahun?
    \item Apa legacy pribadi yang ingin kamu tinggalkan?
    \item Apa pengalaman hidup yang paling ingin kamu miliki?
\end{enumerate}

\subsection*{B. Shared Dreams (Bersama)}
\begin{enumerate}
    \item[4.] Apa yang ingin kita capai \textbf{bersama} dalam 1 tahun?
    \item[5.] Apa yang ingin kita capai \textbf{bersama} dalam 5 tahun?
    \item[6.] Apa yang ingin kita capai \textbf{bersama} dalam 20 tahun?
    \item[7.] Bagaimana kita ingin diingat oleh anak/cucu kita?
\end{enumerate}

\subsection*{C. Alignment Check}
Untuk setiap tujuan di atas, diskusikan:
\begin{itemize}
    \item Apakah tujuan pribadi kita saling mendukung atau bertentangan?
    \item Apa yang perlu dikorbankan untuk tujuan bersama?
    \item Bagaimana kita akan memantau progress?
\end{itemize}

\textbf{Action Items:}
\begin{enumerate}
    \item \rule{12cm}{0.4pt}
    \item \rule{12cm}{0.4pt}
    \item \rule{12cm}{0.4pt}
\end{enumerate}
\end{exercisebox}

\section{Level 4: Symbols---Makna di Balik Objek dan Konsep}

\keyterm{Symbols} adalah objek, tempat, atau konsep yang memiliki makna lebih dalam bagi hubungan. Ini adalah ``bahasa rahasia'' pasangan.

\subsection{Kategori Simbol}

\begin{center}
\renewcommand{\arraystretch}{1.6}
\begin{tabular}{|>{\raggedright}p{3cm}|>{\raggedright}p{4.5cm}|>{\raggedright}p{4.5cm}|>{\raggedright\arraybackslash}p{3cm}|}
\hline
\rowcolor{primary!20}
\textbf{Simbol} & \textbf{Makna Potensial} & \textbf{Diskusi yang Diperlukan} & \textbf{Contoh Konflik} \\
\hline
\textbf{Home} & Security, creativity, status, comfort & Apa fungsi utama rumah? Apa yang membuatnya ``berhasil''? & Minimalis vs collector \\
\hline
\textbf{Money} & Security, freedom, success, enjoyment & Mengapa kita menabak/belanja? Uang untuk apa? & Saver vs spender \\
\hline
\textbf{Time} & Productivity, presence, achievement, connection & Apa makna waktu yang ``well-spent''? & Workaholic vs quality time seeker \\
\hline
\textbf{Sex} & Intimacy, pleasure, obligation, power & Apa makna keintiman fisik? & High vs low desire \\
\hline
\textbf{Work} & Identity, provision, passion, burden & Apa peran kerja dalam hidup? & Ambisi vs balance \\
\hline
\textbf{Family} & Obligation, joy, burden, identity & Seberapa dekat dengan keluarga besar? & Enmeshed vs independent \\
\hline
\end{tabular}
\end{center}

\subsection{Menjelajahi Makna Simbol}

\begin{tipbox}[Dialog tentang Simbol]
\textbf{Home:}
\begin{itemize}
    \item ``Ceritakan tentang rumah masa kecilmu. Bagaimana rasanya?''
    \item ``Apa yang membuatmu merasa `pulang'?''
    \item ``Rumah ideal kita seperti apa? Kenapa?''
\end{itemize}

\textbf{Money:}
\begin{itemize}
    \item ``Apa pengalaman pertamamu dengan uang?''
    \item ``Apa yang membuatmu cemas tentang uang?''
    \item ``Uang mewakili apa bagimu: security, freedom, atau yang lain?''
\end{itemize}

\textbf{Time:}
\begin{itemize}
    \item ``Apa memori terbaikmu tentang waktu bersama keluarga?''
    \item ``Kapan terakhir kali kamu merasa `hidup'? Apa yang terjadi?''
    \item ``Bagaimana kamu ingin menghabiskan waktu kalau hanya punya setahun lagi?''
\end{itemize}
\end{tipbox}

\section{Case Studies: Membangun Shared Meaning dari Nol}

\subsection{Case Study 1: Alex dan Maya---Pasangan Baru}

\begin{praktikbox}[Background]
\textbf{Situation:} Alex dan Maya menikah 2 tahun. Keduanya fokus pada karir masing-masing. Hubungan mulai terasa hampa---seperti ``roommate yang berbagi tempat tidur.''

\textbf{Assessment:}
\begin{itemize}
    \item Rituals: Hampir tidak ada rutinitas bersama
    \item Roles: Belum jelas ekspektasi masing-masing
    \item Goals: Target pribadi yang terpisah
    \item Symbols: Belum membicarakan makna uang, waktu, rumah
\end{itemize}
\end{praktikbox}

\begin{praktikbox}[Intervensi Shared Meaning]
\textbf{Step 1: Ritual Building}
\begin{itemize}
    \item Dibuat: Morning coffee ritual (15 menit tanpa HP)
    \item Dibuat: Weekly ``dream date'' (diskusi tentang impian, bukan tugas)
    \item Dibuat: Monthly ``adventure day'' (mencoba aktivitas baru bersama)
\end{itemize}

\textbf{Step 2: Roles Clarification}
\begin{itemize}
    \item Diskusi panjang tentang makna ``menjadi pasangan''
    \item Kesepakatan: ``Kami adalah tim. Kemenangan satu adalah kemenangan keduanya.''
    \item Pembagian tanggung jawab rumah tangga yang jelas
\end{itemize}

\textbf{Step 3: Goals Alignment}
\begin{itemize}
    \item Ditemukan: Keduanya ingin punya bisnis bersama suatu hari nanti
    \item Dibuat: 5-year plan dengan milestone bersama
    \item Dibuat: Annual goal review tradition
\end{itemize}

\textbf{Step 4: Symbols Exploration}
\begin{itemize}
    \item Alex: Uang = freedom (dari kemiskinan masa kecil)
    \item Maya: Uang = security (orang tuanya bercerai karena utang)
    \item Agreement: 60\% needs, 20\% fun money masing-masing, 20\% dreams
\end{itemize}
\end{praktikbox}

\begin{praktikbox}[Result: After 6 Months]
\begin{itemize}
    \item Morning ritual menjadi ``sacred''---tidak pernah dilewatkan meski sibuk
    \item Monthly adventure date menjadi highlight yang dinantikan
    \item Shared business plan sedang dikerjakan---rasa tim terbentuk kuat
    \item Konflik uang berkurang karena saling memahami makna di baliknya
\end{itemize}
\end{praktikbox}

\subsection{Case Study 2: Budi dan Sari---Pasangan 15 Tahun}

\begin{praktikbox}[Background]
\textbf{Situation:} Budi dan Sari menikah 15 tahun, punya 2 anak remaja. Fokus sepenuhnya pada anak. Setelah anak mulai mandiri, mereka menyadari mereka ``tidak mengenal'' satu sama lain lagi.

\textbf{The Drift:}
\begin{itemize}
    \item Semua percakapan tentang anak atau logistik
    \item Tidak ada ritual berdua sejak anak lahir
    \item Peran: ``Ayah'' dan ``Ibu'' lebih kuat dari ``Suami'' dan ``Istri''
    \item Goals: Sepenuhnya terkait anak---tidak ada untuk mereka berdua
\end{itemize}
\end{praktikbox}

\begin{praktikbox}[Intervensi Shared Meaning]
\textbf{Step 1: Rekindling Rituals}
\begin{itemize}
    \item Dibuat: ``Second honeymoon''---trip berdua pertama dalam 10 tahun
    \item Dibuat: Weekly coffee date (seminggu sekali saat anak di sekolah)
    \item Dibuat: Nightly 10-minute check-in (sebelum tidur, tanpa anak)
\end{itemize}

\textbf{Step 2: Redefining Roles}
\begin{itemize}
    \item Diskusi: ``Siapa kita selain orang tua?''
    \item Discovery: Budi ingin kembali bermusik; Sari ingin menulis
    \item Agreement: Masing-masing punya waktu untuk passion pribadi
\end{itemize}

\textbf{Step 3: New Shared Goals}
\begin{itemize}
    \item Dibuat: ``Empty nest vision''---rencana untuk saat anak pergi
    \item Dibuat: Travel bucket list
    \item Dibuat: Joint hobby (memasak kursus)
\end{itemize}
\end{praktikbox}

\subsection{Case Study 3: David dan Rina---Visi Berbeda}

\begin{warningbox}[The Challenge]
\textbf{David:} Island type. Ingin pensiun dini, hidup sederhana di pegunungan, jauh dari keramaian.

\textbf{Rina:} Wave type. Ingin sukses besar, hidup di kota, aktif dalam komunitas sosial.

\textbf{Konflik:} Perbedaan visi hidup fundamental---bisa menjadi perpetual problem yang tidak terselesaikan.
\end{warningbox}

\begin{praktikbox}[Resolusi: Finding Shared Meaning dalam Perbedaan]
\textbf{Step 1: Dreams Within Conflict Dialogue}
\begin{itemize}
    \item David: Pegunungan = peace, recovery dari trauma kerja, kebebasan
    \item Rina: Kota = connection, purpose, validation, membuat perbedaan
\end{itemize}

\textbf{Step 2: Finding Common Ground}
\begin{itemize}
    \item Ditemukan: Keduanya ingin ``membuat perbedaan''---cara berbeda
    \item Ditemukan: Keduanya butuh ``ruang aman''---definisi berbeda
    \item Ditemukan: Keduanya menghargai ``kebebasan''---ekspresi berbeda
\end{itemize}

\textbf{Step 3: Creative Compromise}
\begin{itemize}
    \item Agreement: 6 bulan kota, 6 bulan pegunungan (setelah pensiun)
    \item Agreement: David akan volunteer di komunitas Rina
    \item Agreement: Rina akan bantu David membangun komunitas di pegunungan
    \item Agreement: Selama ini---liburan regular ke kedua tempat
\end{itemize}
\end{praktikbox}

\section{Ketika Pasangan Punya Visi Berbeda: Cara Negosiasi}

\begin{warningbox}[Realitas Perbedaan]
Tidak semua pasangan akan sepenuhnya setuju tentang shared meaning. Perbedaan dalam:
\begin{itemize}
    \item Nilai fundamental (religious, moral)
    \item Lifestyle preferences (city vs country, simple vs luxurious)
    \item Life goals (children vs childfree, career vs family focus)
    \item Attachment needs (Island vs Wave patterns)
\end{itemize}

\textit{Ini adalah \textbf{perpetual problems} yang perlu dikelola, bukan diselesaikan.}
\end{warningbox}

\subsection{Framework Negosiasi Shared Meaning}

\begin{center}
\begin{tikzpicture}[
    step/.style={draw=#1, line width=2pt, rounded corners=10pt, fill=#1!10,
                  minimum width=12cm, minimum height=1.5cm, align=left, text width=11cm}
]
    \node[step=sagegreen] (s1) at (0,8) {
        \textbf{Step 1: EXPLORASI}\\
        \small Masing-masing menjelaskan dream/nilai tanpa interruption atau defense
    };
    
    \node[step=secondary] (s2) at (0,5.5) {
        \textbf{Step 2: VALIDASI}\\
        \small Masing-masing memvalidasi: ``Aku mengerti ini penting karena...''
    };
    
    \node[step=accent] (s3) at (0,3) {
        \textbf{Step 3: CARI OVERLAP}\\
        \small Identifikasi nilai/goal yang sama, meski ekspresinya berbeda
    };
    
    \node[step=primary] (s4) at (0,0.5) {
        \textbf{Step 4: KREATIF KOMPROMI}\\
        \small Brainstorm solusi yang menghormati keduanya---pikir di luar kotak
    };
    
    \node[step=warmgray] (s5) at (0,-2) {
        \textbf{Step 5: AGREEMENT DAN REVIEW}\\
        \small Buat agreement konkret dan jadwalkan review regular
    };
    
    \draw[->, line width=2pt] (s1) -- (s2);
    \draw[->, line width=2pt] (s2) -- (s3);
    \draw[->, line width=2pt] (s3) -- (s4);
    \draw[->, line width=2pt] (s4) -- (s5);
\end{tikzpicture}
\end{center}

\subsection{Teknik Kompromi Kreatif}

\begin{tipbox}[Strategi Negosiasi]
\textbf{1. The Time-Division Solution}
\begin{itemize}
    \item ``6 bulan di A, 6 bulan di B''
    \item ``Weekdays untuk X, weekends untuk Y''
    \item ``Fase 1: X, Fase 2: Y''
\end{itemize}

\textbf{2. The Hybrid Solution}
\begin{itemize}
    \item Gabungkan elemen dari keduanya
    \item Contoh: Hidup di kota dengan balkon taman (city + nature)
\end{itemize}

\textbf{3. The Autonomy Agreement}
\begin{itemize}
    \item ``Kamu lakukan X, aku lakukan Y, kita respect pilihan masing-masing''
    \item Contoh: Satu fokus karir, satu fokus keluarga---tapi keduanya valid
\end{itemize}

\textbf{4. The Trial Period}
\begin{itemize}
    \item Coba satu cara selama periode tertentu, lalu evaluasi
    \item Coba cara lain, bandingkan
\end{itemize}

\textbf{5. The Third Way}
\begin{itemize}
    \item Brainstorm opsi C yang tidak dipikirkan sebelumnya
    \item Minta input dari counselor, mentor, atau teman bijak
\end{itemize}
\end{tipbox}

\section{Integrasi dengan Attachment Styles}

\begin{center}
\renewcommand{\arraystretch}{1.7}
\begin{tabular}{|>{\raggedright}p{2.5cm}|>{\raggedright}p{5cm}|>{\raggedright}p{5cm}|}
\hline
\rowcolor{primary!20}
\textbf{Attachment} & \textbf{Biasa dalam Shared Meaning} & \textbf{Strategi Optimal} \\
\hline
\textbf{Anchor} & Natural dalam membangun shared meaning; balance individual dan shared goals & Gunakan kekuatan untuk memfasilitasi diskusi; bantu pasangan yang struggle \\
\hline
\textbf{Island} & Cenderung fokus pada goals individual; sulit dengan ``interdependent'' rituals & Beri ruang autonomy dalam shared meaning; ritual jangan terlalu ``melekat'' \\
\hline
\textbf{Wave} & Kuat dalam emotional rituals; mungkin terlalu fokus pada togetherness & Balance dengan individual goals; pastikan tidak mengorbankan diri sendiri \\
\hline
\end{tabular}
\end{center}

\subsection{Challenges by Combination}

\begin{praktikbox}[Island + Wave dalam Shared Meaning]
\textbf{Challenge:}
\begin{itemize}
    \item Wave ingin banyak shared rituals (koneksi)
    \item Island ingin individual space (autonomy)
    \item Risiko: Wave merasa ditinggal, Island merasa tercekik
\end{itemize}

\textbf{Solution:}
\begin{itemize}
    \item Rituals yang ada ``escape clause''-nya
    \item Balance: Some together, some apart
    \item Agreement yang jelas tentang ``me time''
\end{itemize}
\end{praktikbox}

\begin{praktikbox}[Two Islands dalam Shared Meaning]
\textbf{Challenge:}
\begin{itemize}
    \item Keduanya fokus pada goals individual
    \item Risiko: Hubungan menjadi ``parallel lives''
    \item Sedikit emotional investment dalam ``bersama''
\end{itemize}

\textbf{Solution:}
\begin{itemize}
    \item Scheduled shared meaning activities (wajib)
    \item External accountability (counselor, friends)
    \item Practical shared goals (bisnis, proyek rumah)
\end{itemize}
\end{praktikbox}

\begin{praktikbox}[Two Waves dalam Shared Meaning]
\textbf{Challenge:}
\begin{itemize}
    \item Terlalu banyak fokus pada ``kita'', sedikit individual identity
    \item Risiko: Codependency, loss of self
    \item Emosi yang terlalu terkait satu sama lain
\end{itemize}

\textbf{Solution:}
\begin{itemize}
    \item Sengaja maintain individual friendships/hobbies
    \item Rituals yang support individual growth
    \item Regular check: ``Apa kebutuhanku sendiri?''
\end{itemize}
\end{praktikbox}

\section{Annual Relationship Review: Template}

\begin{exercisebox}[TEMPLATE: Annual Relationship Review]
\textbf{Tanggal:} \rule{4cm}{0.4pt} \textbf{Tempat:} \rule{5cm}{0.4pt}

\textit{(Blokir 3-4 jam. Bawa notes tahun lalu jika ada.)}

\subsection*{BAGIAN A: Year in Review (30 menit)}
\begin{enumerate}
    \item Apa pencapaian terbesar hubungan kita tahun ini?
    \item Apa tantangan terbesar yang kita hadapi?
    \item Apa momen paling berkesan bersama?
    \item Apa yang kita pelajari tentang satu sama lain?
    \item Apa yang kita pelajari tentang hubungan kita?
\end{enumerate}

\subsection*{BAGIAN B: Rituals Review (20 menit)}
\begin{enumerate}
    \item Ritual apa yang berjalan baik tahun ini?
    \item Ritual apa yang terbengkalai dan perlu di-revive?
    \item Ritual baru apa yang ingin kita coba tahun depan?
    \item Apa yang membuat ritual kita berharga?
\end{enumerate}

\subsection*{BAGIAN C: Roles Check-in (20 menit)}
\begin{enumerate}
    \item Apakah peran kita sebagai pasangan masih sesuai ekspektasi?
    \item Ada perubahan dalam peran (karir, parenting, dll)?
    \item Apa yang ingin kita ubah dalam cara kita menjalankan peran?
\end{enumerate}

\subsection*{BAGIAN D: Goals Progress (30 menit)}
\begin{enumerate}
    \item Review goals tahun lalu---apa yang tercapai?
    \item Apa yang belum tercapai dan mengapa?
    \item Goals baru untuk tahun depan:
    \begin{itemize}
        \item Individual: \rule{8cm}{0.4pt}
        \item Shared: \rule{8cm}{0.4pt}
    \end{itemize}
\end{enumerate}

\subsection*{BAGIAN E: Symbols Exploration (20 menit)}
\begin{enumerate}
    \item Apakah makna dari rumah, uang, waktu berubah tahun ini?
    \item Apa yang perlu kita klaim ulang atau redefinisikan?
\end{enumerate}

\subsection*{BAGIAN F: Looking Forward (20 menit)}
\begin{enumerate}
    \item Visi besar: Di mana kita ingin berada dalam 5 tahun?
    \item Prioritas #1 untuk tahun depan:
    \item Komitmen masing-masing:
    \begin{itemize}
        \item Partner A: \rule{8cm}{0.4pt}
        \item Partner B: \rule{8cm}{0.4pt}
    \end{itemize}
\end{enumerate}

\subsection*{BAGIAN G: Celebration and Closing (10 menit)}
\begin{enumerate}
    \item Apresiasi untuk pasangan:
    \item Janji untuk tahun depan:
    \item Ritual penutup (pelukan, ciuman, doa bersama, dll):
\end{enumerate}

\vspace{0.5cm}
\textbf{Ditandatangani:}\\
Partner A: \rule{5cm}{0.4pt} Tanggal: \rule{3cm}{0.4pt}\\
Partner B: \rule{5cm}{0.4pt} Tanggal: \rule{3cm}{0.4pt}

\vspace{0.5cm}
\textit{(Simpan ini. Review lagi tahun depan!)}
\end{exercisebox}

\section{Action Steps: Memulai Shared Meaning Journey}

\begin{exercisebox}[Program 30 Hari: Membangun Shared Meaning]

\textbf{MINGGU 1: Assessment dan Rituals}
\begin{itemize}
    \item \textbf{Hari 1-2:} Kerjakan Shared Meaning Assessment bersama
    \item \textbf{Hari 3-4:} Identifikasi ritual yang sudah ada (meski tidak disadari)
    \item \textbf{Hari 5-7:} Buat satu ritual baru dan mulai praktikkan
\end{itemize}

\textbf{MINGGU 2: Roles dan Goals}
\begin{itemize}
    \item \textbf{Hari 8-10:} Diskusikan makna peran masing-masing
    \item \textbf{Hari 11-14:} Kerjakan Goals Alignment worksheet
\end{itemize}

\textbf{MINGGU 3: Symbols dan Stories}
\begin{itemize}
    \item \textbf{Hari 15-17:} Dialog tentang simbol: rumah, uang, waktu
    \item \textbf{Hari 18-21:} Ceritakan ``story of us''---bagaimana hubungan ini berkembang
\end{itemize}

\textbf{MINGGU 4: Integration dan Commitment}
\begin{itemize}
    \item \textbf{Hari 22-25:} Buat Couple Agreement tentang shared meaning
    \item \textbf{Hari 26-28:} Jadwalkan Annual Review pertama
    \item \textbf{Hari 29-30:} Reflection dan celebration
\end{itemize}
\end{exercisebox}

\section{Kesimpulan: The Ultimate Goal}

\begin{insightbox}
\textbf{Key Takeaways Chapter 10:}

\begin{enumerate}
    \item \textbf{Shared meaning adalah puncak dari Seven Principles}---tempat di mana semua prinsip lain bermuara
    
    \item \textbf{Shared meaning melampaui daily logistics}---ini tentang makna hidup yang dibangun bersama
    
    \item \textbf{Empat pilar shared meaning:} Rituals, Roles, Goals, dan Symbols
    
    \item \textbf{Shared meaning bisa dibangun}---meski dimulai dari nol
    
    \item \textbf{Perbedaan dalam visi bukan akhir}---dengan dialog dan kreativitas, kompromi bermakna mungkin
    
    \item \textbf{Annual review menjaga shared meaning tetap hidup}---ini bukan ``set and forget''
\end{enumerate}

\textit{``Hubungan yang berhasil bukan hanya tentang bertahan dari badai. Ini tentang menciptakan peta bersama---menuju destinasi yang kalian bangun berdua.''} --- John Gottman
\end{insightbox}

\begin{tipbox}[Final Checklist]
\begin{itemize}
    \item[$\square$] Shared Meaning Assessment sudah dikerjakan
    \item[$\square$] Minimal satu ritual baru sudah dimulai
    \item[$\square$] Dialog tentang roles sudah dilakukan
    \item[$\square$] Goals alignment sudah dibahas
    \item[$\square$] Makna simbol kunci sudah dieksplorasi
    \item[$\square$] Annual Review sudah dijadwalkan
    \item[$\square$] Couple Agreement tentang shared meaning sudah dibuat
\end{itemize}
\end{tipbox}

\vspace{1cm}
\begin{center}
\begin{tikzpicture}
    \node[draw=primary, line width=3pt, rounded corners=20pt, 
          fill=lightbg, inner sep=30pt, align=center] {
        {\Large\bfseries\color{primary} Selamat!}\\[0.5cm]
        {\large\color{darktext} Kamu telah menyelesaikan semua}\\
        {\large\color{darktext} Seven Principles for Making Marriage Work.}\\[0.5cm]
        {\color{warmgray} Ingat: Ini bukan akhir, tapi awal.}\\
        {\color{warmgray} Shared meaning adalah perjalanan seumur hidup.}
    };
\end{tikzpicture}
\end{center}
